\documentclass[12pt,a4paper]{article}

\usepackage{fontspec}
\usepackage{listings}
\usepackage{xcolor}
\usepackage{graphicx}
\usepackage{geometry}
\usepackage{hyperref}
\usepackage{xepersian}

\geometry{margin=2.5cm}
\settextfont{XB Niloofar}
\setlatintextfont{Times New Roman}

\definecolor{codegreen}{rgb}{0,0.6,0}
\definecolor{codegray}{rgb}{0.5,0.5,0.5}
\definecolor{codepurple}{rgb}{0.58,0,0.82}
\definecolor{backcolour}{rgb}{0.95,0.95,0.92}

\lstdefinestyle{mystyle}{
    backgroundcolor=\color{backcolour},   
    commentstyle=\color{codegreen},
    keywordstyle=\color{magenta},
    numberstyle=\tiny\color{codegray},
    stringstyle=\color{codepurple},
    basicstyle=\ttfamily\footnotesize,
    breakatwhitespace=false,         
    breaklines=true,                 
    captionpos=b,                    
    keepspaces=true,                 
    numbers=left,                    
    numbersep=5pt,                  
    showspaces=false,                
    showstringspaces=false,
    showtabs=false,                  
    tabsize=2
}
\lstset{style=mystyle}

\title{\textbf{گزارش پیاده‌سازی زمان‌بند CFS و COW Fork در xv6}}
\author{دلارام روحانی - ۴۰۲۳۱۰۲۶ \\
امیرعباس انتظاری - ۴۰۲۳۱۰۰۶}
\date{\today}

\begin{document}

\maketitle

\tableofcontents
\newpage

%%%%%%%%%%%%%%%%%%%%%%%%%%%%%%%%%%%%%%%%%%%%%%%%%%%%%%%%%%%%%%%
% Part 1: CFS Scheduler
%%%%%%%%%%%%%%%%%%%%%%%%%%%%%%%%%%%%%%%%%%%%%%%%%%%%%%%%%%%%%%%

\part{پیاده‌سازی زمان‌بند CFS}

\section{مقدمه به CFS}
\textbf{\lr{Completely Fair Scheduler (CFS)}} یک الگوریتم زمان‌بندی است که در هسته لینوکس استفاده می‌شود. هدف اصلی CFS توزیع عادلانه زمان CPU بین تمام پردازه‌ها است.

در xv6 اصلی، از یک الگوریتم ساده \lr{Round-Robin} استفاده می‌شود که هر پردازه را به نوبت اجرا می‌کند. این روش ساده است اما عادلانه نیست زیرا اولویت پردازه‌ها را در نظر نمی‌گیرد.

\subsection{مفاهیم کلیدی CFS}
\begin{itemize}
    \item \textbf{\lr{vruntime (Virtual Runtime)}}: زمان مجازی اجرای هر پردازه. پردازه‌ای که کمتر اجرا شده، \lr{vruntime} کمتری دارد.
    \item \textbf{\lr{Weight}}: وزن هر پردازه بر اساس اولویت. پردازه‌های مهم‌تر وزن بیشتری دارند.
    \item \textbf{\lr{Min-Heap}}: برای انتخاب سریع پردازه با کمترین \lr{vruntime} از یک \lr{Min-Heap} استفاده می‌شود.
\end{itemize}

\section{اهداف پیاده‌سازی CFS}
\begin{itemize}
    \item جایگزینی زمان‌بند \lr{Round-Robin} با CFS
    \item پیاده‌سازی ساختار داده \lr{Min-Heap} برای مدیریت صف اجرا
    \item اضافه کردن فیلدهای \lr{vruntime} و \lr{weight} به ساختار پردازه
    \item به‌روزرسانی \lr{vruntime} در هر بار \lr{yield}
\end{itemize}

\section{جزئیات پیاده‌سازی CFS}

\subsection{تغییرات در ساختار پردازه (proc.h)}
دو فیلد جدید به ساختار \lr{proc} اضافه شد:

\begin{latin}
\begin{lstlisting}[language=C, caption={فیلدهای جدید در struct proc}]
struct proc {
  // ... existing fields ...
  
  int weight;        // Process weight (based on priority)
  uint64 vruntime;   // Virtual runtime for CFS
};
\end{lstlisting}
\end{latin}

\begin{itemize}
    \item \lr{weight}: وزن پردازه. مقدار پیش‌فرض ۱۰۲۴ (مطابق با \lr{nice=0} در لینوکس).
    \item \lr{vruntime}: زمان مجازی اجرا. پردازه با کمترین \lr{vruntime} اولویت بالاتری برای اجرا دارد.
\end{itemize}

\subsection{ساختار \lr{Min-Heap (minheap.h و minheap.c)}}
یک \lr{Min-Heap} برای نگهداری پردازه‌های آماده اجرا پیاده‌سازی شد:

\begin{latin}
\begin{lstlisting}[language=C, caption={ساختار Min-Heap}]
struct minheap {
  struct proc *procs[NPROC];  // Array of process pointers
  int size;                   // Current number of elements
  int capacity;               // Maximum capacity (NPROC)
};
\end{lstlisting}
\end{latin}

توابع اصلی:
\begin{itemize}
    \item \lr{minheap\_init()}: مقداردهی اولیه heap
    \item \lr{minheap\_insert()}: اضافه کردن پردازه به heap
    \item \lr{minheap\_extract\_min()}: استخراج پردازه با کمترین \lr{vruntime}
\end{itemize}

\begin{latin}
\begin{lstlisting}[language=C, caption={تابع heapify\_up}]
static void heapify_up(struct minheap *heap, int idx) {
  while (idx > 0) {
    int p = parent(idx);
    if (heap->procs[idx]->vruntime < heap->procs[p]->vruntime) {
      swap(heap, idx, p);
      idx = p;
    } else {
      break;
    }
  }
}
\end{lstlisting}
\end{latin}

\begin{latin}
\begin{lstlisting}[language=C, caption={تابع extract\_min}]
struct proc* minheap_extract_min(struct minheap *heap) {
  if (heap->size == 0) return 0;
  
  struct proc *min_proc = heap->procs[0];
  
  // Move last element to root
  heap->procs[0] = heap->procs[heap->size - 1];
  heap->size--;
  
  // Heapify down from root
  if (heap->size > 0) {
    heapify_down(heap, 0);
  }
  
  return min_proc;
}
\end{lstlisting}
\end{latin}

\subsection{متغیرهای سراسری \lr{CFS (proc.c)}}
در فایل \lr{proc.c} متغیرهای سراسری زیر اضافه شد:

\begin{latin}
\begin{lstlisting}[language=C, caption={متغیرهای سراسری CFS}]
// CFS scheduler data structures
struct minheap run_queue;     // Min-heap of runnable processes
struct spinlock runq_lock;    // Lock for the run queue
uint64 min_vruntime = 0;      // Global minimum vruntime
\end{lstlisting}
\end{latin}

\subsection{مقداردهی اولیه (procinit)}
در تابع \lr{procinit()} قفل و heap مقداردهی اولیه می‌شوند:

\begin{latin}
\begin{lstlisting}[language=C, caption={مقداردهی اولیه CFS}]
void procinit(void) {
  initlock(&runq_lock, "runqueue");
  minheap_init(&run_queue);
  min_vruntime = 0;
  // ...
}
\end{lstlisting}
\end{latin}

\subsection{تخصیص پردازه (allocproc)}
هنگام ایجاد پردازه جدید، فیلدهای CFS مقداردهی می‌شوند:

\begin{latin}
\begin{lstlisting}[language=C, caption={مقداردهی فیلدهای CFS در allocproc}]
static struct proc* allocproc(void) {
  // ...
  
  // Initialize CFS fields
  p->vruntime = min_vruntime;  // Start with current min
  p->weight = 1024;            // Default weight (nice 0)
  
  return p;
}
\end{lstlisting}
\end{latin}

\subsection{زمان‌بند جدید (scheduler)}
تابع \lr{scheduler()} بازنویسی شد تا از CFS استفاده کند:

\begin{latin}
\begin{lstlisting}[language=C, caption={زمان‌بند CFS}]
void scheduler(void) {
  struct proc *p;
  struct cpu *c = mycpu();
  c->proc = 0;
  
  for(;;){
    intr_on();
    intr_off();

    // CFS: Extract process with minimum vruntime
    acquire(&runq_lock);
    p = minheap_extract_min(&run_queue);
    release(&runq_lock);

    if(p != 0) {
      acquire(&p->lock);
      if(p->state == RUNNABLE) {
        p->state = RUNNING;
        c->proc = p;
        swtch(&c->context, &p->context);
        c->proc = 0;
      }
      release(&p->lock);
    } else {
      asm volatile("wfi");  // Wait for interrupt
    }
  }
}
\end{lstlisting}
\end{latin}

\subsection{به‌روزرسانی \lr{vruntime (yield)}}
در تابع \lr{yield()}، \lr{vruntime} پردازه به‌روز می‌شود:

\begin{latin}
\begin{lstlisting}[language=C, caption={به‌روزرسانی vruntime در yield}]
void yield(void) {
  struct proc *p = myproc();
  acquire(&p->lock);
  
  // CFS: Update vruntime based on weight
  // vruntime += delta * (NICE_0_WEIGHT / weight)
  p->vruntime += (1024 / p->weight);
  
  // Update global min_vruntime
  acquire(&runq_lock);
  if(p->vruntime > min_vruntime)
    min_vruntime = p->vruntime;
  release(&runq_lock);
  
  p->state = RUNNABLE;
  
  // Add back to run queue
  acquire(&runq_lock);
  minheap_insert(&run_queue, p);
  release(&runq_lock);
  
  sched();
  release(&p->lock);
}
\end{lstlisting}
\end{latin}

\subsection{اضافه کردن پردازه به صف \lr{(fork و wakeup)}}
در توابع \lr{fork()} و \lr{wakeup()}، پردازه‌ها به \lr{run\_queue} اضافه می‌شوند:

\begin{latin}
\begin{lstlisting}[language=C, caption={اضافه کردن به صف اجرا}]
// In fork:
np->state = RUNNABLE;
acquire(&runq_lock);
minheap_insert(&run_queue, np);
release(&runq_lock);
\end{lstlisting}
\end{latin}

\section{فرمول محاسبه vruntime}
فرمول به‌روزرسانی \lr{vruntime} به صورت زیر است:

\begin{equation}
vruntime_{new} = vruntime_{old} + \frac{NICE\_0\_WEIGHT}{weight} \times \Delta t
\end{equation}

که در آن:
\begin{itemize}
    \item \lr{NICE\_0\_WEIGHT = 1024}: وزن پیش‌فرض
    \item \lr{weight}: وزن پردازه فعلی
    \item $\Delta t$: زمان اجرا (در پیاده‌سازی ساده برابر ۱)
\end{itemize}

\section{نتیجه‌گیری CFS}
با پیاده‌سازی CFS:
\begin{itemize}
    \item پردازه‌هایی که کمتر اجرا شده‌اند، اولویت بالاتری می‌گیرند
    \item پردازه‌های با وزن بالاتر، \lr{vruntime} کمتری افزایش می‌یابد (زمان بیشتری می‌گیرند)
    \item انتخاب پردازه بعدی با پیچیدگی $O(\log n)$ انجام می‌شود
\end{itemize}

\newpage

%%%%%%%%%%%%%%%%%%%%%%%%%%%%%%%%%%%%%%%%%%%%%%%%%%%%%%%%%%%%%%%
% Part 2: COW Fork
%%%%%%%%%%%%%%%%%%%%%%%%%%%%%%%%%%%%%%%%%%%%%%%%%%%%%%%%%%%%%%%

\part{پیاده‌سازی \lr{COW Fork}}

\section{مقدمه به COW}
در سیستم‌عامل، زمانی که فراخوانی سیستمی \lr{fork()} اجرا می‌شود، یک کپی کامل از فضای آدرس پردازه والد برای پردازه فرزند ایجاد می‌شود. این روش سنتی باعث اتلاف حافظه و زمان می‌شود، به‌خصوص زمانی که فرزند بلافاصله \lr{exec()} را فراخوانی کند.

\textbf{\lr{Copy-on-Write (COW)}} یک تکنیک بهینه‌سازی است که در آن صفحات حافظه بین والد و فرزند به اشتراک گذاشته می‌شوند تا زمانی که یکی از آن‌ها بخواهد بنویسد. در آن لحظه، یک کپی از صفحه ایجاد می‌شود.

\section{اهداف پیاده‌سازی COW}
\begin{itemize}
    \item پیاده‌سازی \lr{cowfork()} که صفحات را به جای کپی کردن، به اشتراک می‌گذارد
    \item مدیریت \lr{Page Fault} برای تشخیص نوشتن روی صفحات COW
    \item پیاده‌سازی شمارش ارجاع (\lr{Reference Counting}) برای صفحات فیزیکی
    \item اضافه کردن سیستم‌کال \lr{physaddr()} برای تست صحت پیاده‌سازی
\end{itemize}

\section{جزئیات پیاده‌سازی}

\subsection{اضافه کردن بیت COW به PTE}
برای تشخیص صفحات COW از صفحات فقط‌خواندنی معمولی، یک بیت جدید به فلگ‌های PTE اضافه کردیم. از بیت \lr{RSW (Reserved for Software)} استفاده شد:

\begin{latin}
\begin{lstlisting}[language=C, caption={تعریف بیت COW در riscv.h}]
#define PTE_V (1L << 0) // valid
#define PTE_R (1L << 1)
#define PTE_W (1L << 2)
#define PTE_X (1L << 3)
#define PTE_U (1L << 4) // user can access
#define PTE_COW (1L << 8) // copy-on-write page (RSW bit)
\end{lstlisting}
\end{latin}

\subsection{پیاده‌سازی شمارش ارجاع در kalloc.c}
برای مدیریت صحیح حافظه مشترک، یک آرایه برای نگهداری تعداد ارجاعات هر صفحه فیزیکی اضافه شد:

\begin{latin}
\begin{lstlisting}[language=C, caption={ساختار شمارش ارجاع}]
#define MAXPAGES 32768
#define PA2IDX(pa) (((uint64)(pa) - KERNBASE) / PGSIZE)

struct {
  struct spinlock lock;
  int count[MAXPAGES];
} pageref;
\end{lstlisting}
\end{latin}

سه تابع برای مدیریت ارجاعات اضافه شد:
\begin{itemize}
    \item \lr{kref\_incr(pa)}: افزایش شمارنده ارجاع
    \item \lr{kref\_decr(pa)}: کاهش شمارنده ارجاع
    \item \lr{kref\_get(pa)}: خواندن مقدار شمارنده
\end{itemize}

همچنین تابع \lr{kfree()} اصلاح شد تا فقط زمانی صفحه را آزاد کند که شمارنده به صفر برسد:

\begin{latin}
\begin{lstlisting}[language=C, caption={تابع kfree اصلاح‌شده}]
void kfree(void *pa) {
  // Only free if reference count reaches 0
  acquire(&pageref.lock);
  if(pageref.count[PA2IDX(pa)] > 1) {
    pageref.count[PA2IDX(pa)]--;
    release(&pageref.lock);
    return;  // Don't free, still referenced
  }
  pageref.count[PA2IDX(pa)] = 0;
  release(&pageref.lock);
  // ... rest of kfree
}
\end{lstlisting}
\end{latin}

\subsection{پیاده‌سازی uvmcopy\_cow در vm.c}
این تابع جایگزین \lr{uvmcopy} در \lr{cowfork} می‌شود. به جای کپی کردن صفحات:
\begin{enumerate}
    \item صفحات را فقط‌خواندنی می‌کند
    \item بیت COW را ست می‌کند
    \item هر دو جدول صفحه (والد و فرزند) را به یک صفحه فیزیکی نگاشت می‌کند
    \item شمارنده ارجاع را افزایش می‌دهد
\end{enumerate}

\begin{latin}
\begin{lstlisting}[language=C, caption={تابع uvmcopy\_cow}]
int uvmcopy_cow(pagetable_t old, pagetable_t new, uint64 sz) {
  for(i = 0; i < sz; i += PGSIZE){
    // ... get PTE and physical address
    
    // If writable, mark as COW and remove write permission
    if(flags & PTE_W) {
      flags = (flags & ~PTE_W) | PTE_COW;
      *pte = PA2PTE(pa) | flags | PTE_V;  // Update parent
    }
    
    // Map same physical page in child
    mappages(new, i, PGSIZE, pa, flags);
    
    // Increment reference count
    kref_incr((void*)pa);
  }
  return 0;
}
\end{lstlisting}
\end{latin}

\subsection{مدیریت \lr{Page Fault} در \lr{trap.c}}
وقتی پردازه‌ای روی صفحه COW می‌نویسد، یک \lr{Store Page Fault} با \lr{scause=15} رخ می‌دهد. در \lr{usertrap()} این را مدیریت می‌کنیم:

\begin{latin}
\begin{lstlisting}[language=C, caption={مدیریت Page Fault در usertrap}]
} else if(r_scause() == 15) {
  uint64 va = r_stval();
  pte_t *pte = walk(p->pagetable, va, 0);
  
  // Check if it's a COW page fault
  if(pte != 0 && (*pte & PTE_V) && (*pte & PTE_COW)) {
    if(cowfault(p->pagetable, va) < 0) {
      setkilled(p);
    }
  } else if(vmfault(p->pagetable, va, 0) == 0) {
    setkilled(p);
  }
}
\end{lstlisting}
\end{latin}

\subsection{تابع cowfault}
این تابع مسئول رسیدگی به خطای COW است:

\begin{latin}
\begin{lstlisting}[language=C, caption={تابع cowfault}]
int cowfault(pagetable_t pagetable, uint64 va) {
  va = PGROUNDDOWN(va);
  pte_t *pte = walk(pagetable, va, 0);
  
  // Check if COW page
  if((*pte & PTE_COW) == 0) return -1;
  
  uint64 pa = PTE2PA(*pte);
  
  // If only one reference, just restore write permission
  if(kref_get((void*)pa) == 1) {
    *pte = (*pte | PTE_W) & ~PTE_COW;
    return 0;
  }
  
  // Allocate new page and copy content
  char *mem = kalloc();
  memmove(mem, (char*)pa, PGSIZE);
  
  // Update PTE with new page, restore write, clear COW
  uint flags = (PTE_FLAGS(*pte) | PTE_W) & ~PTE_COW;
  *pte = PA2PTE(mem) | flags | PTE_V;
  
  // Decrement reference on old page
  kfree((void*)pa);
  
  return 0;
}
\end{lstlisting}
\end{latin}

\subsection{سیستم‌کال cowfork}
تابع \lr{kcowfork()} در \lr{proc.c} مشابه \lr{kfork()} است، با این تفاوت که از \lr{uvmcopy\_cow()} استفاده می‌کند:

\begin{latin}
\begin{lstlisting}[language=C, caption={تابع kcowfork}]
int kcowfork(void) {
  // Allocate process
  if((np = allocproc()) == 0) return -1;

  // Share user memory using COW (not copy!)
  if(uvmcopy_cow(p->pagetable, np->pagetable, p->sz) < 0){
    freeproc(np);
    return -1;
  }
  // ... rest same as kfork
}
\end{lstlisting}
\end{latin}

\subsection{سیستم‌کال physaddr}
برای تست صحت پیاده‌سازی، این سیستم‌کال شماره صفحه فیزیکی یک آدرس مجازی را برمی‌گرداند:

\begin{latin}
\begin{lstlisting}[language=C, caption={سیستم‌کال physaddr}]
uint64 sys_physaddr(void) {
  struct proc *p = myproc();
  uint64 va;
  argaddr(0, &va);
  
  va = PGROUNDDOWN(va);
  pte_t *pte = walk(p->pagetable, va, 0);
  uint64 pa = PTE2PA(*pte);
  
  return pa / PGSIZE;  // Return page number
}
\end{lstlisting}
\end{latin}

\section{ثبت سیستم‌کال‌ها}
دو سیستم‌کال جدید به فایل‌های زیر اضافه شد:
\begin{itemize}
    \item \lr{syscall.h}: تعریف \lr{SYS\_cowfork=24} و \lr{SYS\_physaddr=25}
    \item \lr{syscall.c}: ثبت توابع handler
    \item \lr{usys.pl}: تولید stub‌های سمت کاربر
    \item \lr{user.h}: اعلان توابع برای برنامه‌های کاربر
\end{itemize}

\section{برنامه تست}
برنامه \lr{cowtest.c} صحت پیاده‌سازی را بررسی می‌کند:

\begin{latin}
\begin{lstlisting}[language=C, caption={برنامه تست COW}]
int main() {
  char *buf = malloc(4096);
  buf[0] = 'A';
  
  int parent_page_before = physaddr(buf);
  
  if(cowfork() == 0) {
    // Child
    int child_page_before = physaddr(buf);
    buf[0] = 'C';  // Triggers COW!
    int child_page_after = physaddr(buf);
    
    // child_page_before == child_page_after means COW failed
    // child_page_before != child_page_after means COW worked
  }
  
  wait(0);
  // Parent's buf[0] should still be 'A'
}
\end{lstlisting}
\end{latin}

\section{نتیجه‌گیری COW}
پیاده‌سازی \lr{COW Fork} با موفقیت انجام شد. خروجی مورد انتظار:
\begin{itemize}
    \item قبل از نوشتن: شماره صفحه والد و فرزند یکسان است (اشتراک‌گذاری)
    \item بعد از نوشتن توسط فرزند: شماره صفحه فرزند تغییر می‌کند (کپی شده)
    \item داده والد تغییر نمی‌کند (ایزوله‌سازی صحیح)
\end{itemize}

این پیاده‌سازی باعث صرفه‌جویی قابل توجه در حافظه و زمان می‌شود، به‌خصوص در سناریوهای \lr{fork+exec}.

\end{document}
